\documentclass[12pt,a4paper]{article}
\usepackage[margin=2.54cm]{geometry}
\usepackage{amsmath,amssymb}
\usepackage{times}
\usepackage[utf8]{inputenc}
\usepackage[T1]{fontenc}
\usepackage{textcomp}
\usepackage{booktabs}
\usepackage{array}

\title{\textbf{SOLVING THE THREE OPEN PROBLEMS OF BLACK HOLE PROPULSION}}
\author{B.\ Greff \\ Independent Researcher \\ Correspondence: ben@thegreffs.com}
\date{}

\begin{document}
\maketitle
\thispagestyle{empty}

\begin{abstract}
\noindent Black hole propulsion --- using Hawking radiation from a micro black hole as a drive --- has been stalled by three unsolved problems: formation (gamma-ray kugelblitz formation are impossible [1]), energy capture (no material reflects GeV radiation, limiting velocity to 0.0001c [2]) and confinement (no mechanism proposed [3]). Solutions to all three are presented under six assumptions. The black hole is formed by gravitational collapse of BPS monopoles and grown to operating mass by railgun feeding. A color-flavor-locked strange quark matter shell absorbs the full Hawking spectrum and re-emits it as magnetically directable keV pairs, sidestepping the reflection problem. An anti-Helmholtz magnetic bottle confines the magnetically charged black hole and funnels pair exhaust through a tapered nozzle. The architecture offers three properties no other concept shares: continuously adjustable thrust (power scales as 1/M$^{2}$), any-matter fuel (conversion at mc$^{2}$ regardless of composition) and self-regulating thermal management (the shell's Usov equilibrium). At 10$^{9}$ kg the system produces 55 MN and reaches Alpha Centauri in 79 years with all fuel carried from launch; at 5 x 10$^{7}$ kg, 0.45c peak velocity yields 11-year transit. Each assumption would constitute a major discovery; the conjunction is speculative but self-consistent and no known physics is violated.
\end{abstract}
\noindent \textit{Keywords: Hawking radiation, interstellar propulsion, strange quark matter, magnetic monopoles, black hole thermodynamics}
\bigskip
\section{INTRODUCTION}
Interstellar propulsion is constrained by a fundamental tension between specific impulse and thrust. Chemical rockets deliver exhaust velocities of order 4 km/s, four orders of magnitude below interstellar requirements. Electric propulsion extends specific impulse but at millinewton thrust densities. A photon rocket requires petawatts per meganewton of thrust and no known energy source meets this requirement. Nuclear fission yields $\sim$10$^{-3}$ mc$^{2}$ per kilogram; fusion $\sim$10$^{-2}$ mc$^{2}$; matter-antimatter annihilation achieves mc$^{2}$ but requires antimatter production many orders of magnitude beyond foreseeable capability [4]. The interstellar propulsion problem is a power-source problem.

\subsection{Hawking radiation as a power source}
Crane and Westmoreland [3] identified a candidate solution: a micro black hole as a Hawking-radiation engine. A black hole of mass M radiates at the Hawking temperature $T_H$ = $\hbar$c$^{3}$/(8$\pi$GM$k_B$), with total power P = $\hbar$c$^{6}$f/(G$^{2}$M$^{2}$) [5, 6]. At M $\sim$ 10$^{9}$ kg the Hawking temperature is kT $\sim$ 10 GeV and total power is approximately 60,000 TW. The concept is elegant: the black hole converts its own rest mass to radiation with near-unit efficiency (the only irreducible loss is the $\sim$17\% neutrino fraction), the fuel is any matter (the spectrum depends on mass and charge, not composition) and the mass-loss rate is self-regulating (feeding grows the hole, underfeeding shrinks it). The black hole is a universal mass-to-radiation converter with a built-in throttle.

\subsection{The three open problems}
Despite this promise, three unsolved problems have stalled the concept for over fifteen years.

Problem 1 --- Formation. Crane and Westmoreland proposed gamma-ray laser implosion (a kugelblitz). Alvarez-Dominguez et al. [1] proved this impossible: Schwinger pair production dissipates the electromagnetic field before gravitational collapse. Page [7] showed idealized photon configurations can in principle form black holes but conceded practical impossibility.

Problem 2 --- Energy capture. The spectrum at kT $\sim$ 10 GeV is dominated by quarks, gluons and leptons at energies where no known material reflects. Lee [2] modeled titanium absorption at 33 km standoff, computing a capture fraction of 46.5\% and maximum velocity of 0.0001c --- functionally useless. Without full-spectrum capture, the black hole is merely a bright isotropic source.

Problem 3 --- Confinement. The black hole must remain centered within the capture geometry while the ship accelerates. Crane and Westmoreland suggested "particle beams"; no detailed mechanism or stability analysis has appeared.

\subsection{This paper}
Solutions to all three problems are presented. The black hole is formed by gravitational collapse of BPS magnetic monopoles --- topological defects whose mutual magnetic repulsion and scalar attraction cancel exactly in the Bogomolny-Prasad-Sommerfield limit [8, 9], leaving gravity as the only long-range force. A spherical shell of color-flavor-locked (CFL) strange quark matter absorbs the full GeV Hawking spectrum via nuclear-density Bethe-Bloch stopping and re-emits it as magnetically directable keV electron-positron pairs through the Usov electrosphere mechanism [10] --- a thermal wavelength converter that sidesteps the reflection problem entirely. An anti-Helmholtz magnetic bottle, acting on the black hole's inherited magnetic charge, confines it near the geometric center and funnels pair exhaust through a tapered nozzle.

The resulting architecture has three properties no other propulsion concept offers simultaneously: thrust is continuously adjustable over orders of magnitude (power scales as 1/M$^{2}$ and the black hole mass is a free design parameter); the fuel is any matter (the black hole converts mass to radiation at mc$^{2}$ per kilogram regardless of composition); and thermal management is solved by design (the SQM shell's Usov equilibrium is self-regulating, with no thermal scaling limit on thrust).

The design rests on six physics assumptions (Table 1). Each would constitute a major discovery if confirmed; the conjunction is speculative. The design is self-consistent: every numerical claim is derived from first principles and anchored to published literature. The point is not that the design is feasible today --- it is not --- but that nothing in the physics chain is known to fail.


\begin{table}[h]
\centering
\caption{Assumption budget. All six assumptions are required simultaneously.}
\label{tab:assumptions}
\small
\begin{tabular}{clll}
\hline
\textbf{\#} & \textbf{Assumption} & \textbf{Status} & \textbf{If wrong} \\
\hline
1 & \parbox[t]{4.5cm}{Hawking radiation exists and is semiclassical at $M \sim 10^9$\,kg} & \parbox[t]{4cm}{Near-universal consensus; never observed} & No power source \\[6pt]
2 & GUT magnetic monopoles exist & \parbox[t]{4cm}{Predicted by all GUTs; never observed [11]} & No formation mechanism \\[6pt]
3 & \parbox[t]{4.5cm}{BPS limit ($\lambda = 0$ at GUT scale)} & \parbox[t]{4cm}{Protected by $N{=}2$ SUSY [12]; untested} & \parbox[t]{3cm}{Monopole cloud cannot collapse} \\[6pt]
4 & \parbox[t]{4.5cm}{Bodmer--Witten hypothesis (SQM is ground state)} & \parbox[t]{4cm}{Bai and Chen 2025 constrains unpaired SQM but not CFL} & No shell material \\[6pt]
5 & \parbox[t]{4.5cm}{CFL is ground state at operating chemical potential} & \parbox[t]{4cm}{Theoretical [13, 14]; never observed} & \parbox[t]{3cm}{Gap-decoupling firewall fails} \\[6pt]
6 & \parbox[t]{4.5cm}{Core-overlap monopole breeding at $O(1)$ rate} & \parbox[t]{4cm}{Derived from duality [12]; unsimulated} & \parbox[t]{3cm}{Cannot manufacture monopole inventory} \\
\hline
\end{tabular}
\end{table}

\subsection{Response to Bai and Chen 2025}
Bai and Chen [15] derive a QCD vacuum energy implying that fully deconfined quark matter has energy per baryon $\sim$1130---1320 MeV --- heavier than the 930 MeV iron threshold and therefore unstable. As they state, this applies "only to this specific scenario and does not extend to quark nuggets with partially restored confinement and chiral symmetry." The CFL phase is precisely characterized by a diquark condensate that partially restores chiral symmetry through a distinct mechanism [13]. Their result constrains unpaired strange quark matter but does not directly constrain CFL stability, which remains an open question for lattice QCD.

\section{BLACK HOLE PARAMETERS AND HAWKING SPECTRUM}
\subsection{Design point}
The reference design adopts M = 10$^{9}$ kg. This is one point in a continuous design space: M is a free parameter and the architecture scales to any mass with only shell-thickness adjustment. The reference mass sits at the high-thrust end of the Crane-Westmoreland viable range [3], chosen for its short-mission capability --- Alpha Centauri in 79 years with all fuel carried from launch --- at the cost of an unfed evaporation lifetime of only 15.8 years.

\subsection{Key parameters}
Hawking temperature:

\begin{equation}
kT = \frac{\hbar c^3}{8\pi G M} = 10.57 \text{ GeV}
\end{equation}
Schwarzschild radius:

\begin{equation}
r_s = \frac{2GM}{c^2} = 1.485 \times 10^{-18} \text{ m}
\end{equation}
Total radiated power:

\begin{equation}
P = \frac{\hbar c^6 f}{G^2 M^2} = 60{,}200 \text{ TW}
\end{equation}
where f = 3.503 x 10$^{-3}$.

Unfed evaporation lifetime:

\begin{equation}
\tau = \frac{G^2 M^3}{3\hbar c^4 f} = 15.8 \text{ yr}
\end{equation}
At this mass the Schwarzschild radius is roughly a thousandth of a proton radius. The Hawking temperature of 10.57 GeV places the thermal peak where the black hole copiously emits all Standard Model particles lighter than about 10 GeV: all six quark flavors except the top ($m_t$/kT = 16.4), all charged leptons, all neutrinos, gluons, photons and --- marginally --- the W, Z and Higgs bosons. The mass-loss rate is dM/dt = P/c$^{2}$ = 670 g/s = 21,140 tonnes per year --- the fueling rate required to maintain the black hole at constant mass.

\subsection{The emission factor f}
The emission factor sums greybody-weighted power contributions from all Standard Model species [6, 16]:

\begin{equation}
f = \sum_i g_i \, \alpha_{s_i} \, S(m_i/kT)
\end{equation}
where $g_i$ counts degrees of freedom, $\alpha_{s_i}$ is the Page power-rate greybody coefficient for spin $s_i$ ($\alpha_{1/2}$ = 4.09 x 10$^{-5}$ per Weyl fermion, $\alpha_1$ = 1.68 x 10$^{-5}$ per vector polarization, $\alpha_0$ = 7.24 x 10$^{-5}$ for scalars, $\alpha_2$ = 1.92 x 10$^{-6}$ per graviton polarization) and S($m_i$/kT) is the thermal-integral suppression factor. The commonly used Boltzmann approximation exp(-m/kT) over-suppresses species with m/kT of order unity by 30---50\%; the correct suppression uses the full Fermi-Dirac thermal integral.

The light quarks and charm dominate (56\% of total power). The primary neutrino fraction is 7.0\% and gravitons add 0.1\%, for a total primary loss of 7.1\%. Species contributions to f are detailed in Table 2.


\begin{table}[h]
\centering
\caption{Species contributions to $f$.}
\label{tab:species}
\small
\begin{tabular}{lllll}
\hline
\textbf{Species} & \textbf{DOF} & \textbf{$m/kT$} & \textbf{$S(m/kT)$} & \textbf{Contribution to $f$} \\
\hline
$u, d, s$ quarks & 36 Weyl & $\sim$0 & 1.000 & $1.472 \times 10^{-3}$ \\
$c$ quark & 12 Weyl & 0.120 & 1.000 & $4.91 \times 10^{-4}$ \\
$b$ quark & 12 Weyl & 0.395 & 0.9995 & $4.91 \times 10^{-4}$ \\
$t$ quark & 12 Weyl & 16.36 & $1.2 \times 10^{-5}$ & $3.6 \times 10^{-8}$ \\
$e, \mu$ leptons & 8 Weyl & $\sim$0 & 1.000 & $3.27 \times 10^{-4}$ \\
$\tau$ lepton & 4 Weyl & 0.168 & 1.000 & $1.64 \times 10^{-4}$ \\
Neutrinos (3 gen.) & 6 Weyl & 0 & 1.000 & $2.45 \times 10^{-4}$ \\
Gluons & 16 pol. & 0 & 1.000 & $2.69 \times 10^{-4}$ \\
Photon & 2 pol. & 0 & 1.000 & $3.36 \times 10^{-5}$ \\
$W^{\pm}$ & 6 pol. & 7.60 & 0.0511 & $5.15 \times 10^{-6}$ \\
$Z$ & 3 pol. & 8.63 & 0.0255 & $1.29 \times 10^{-6}$ \\
Higgs & 1 & 11.83 & 0.0024 & $1.74 \times 10^{-7}$ \\
Graviton & 2 pol. & 0 & 1.000 & $3.84 \times 10^{-6}$ \\
\hline
TOTAL & & & & $f = 3.503 \times 10^{-3}$ \\
\hline
\end{tabular}
\end{table}

\subsection{Magnetic charge and sub-extremality}
The black hole carries net magnetic charge g = 1.182 x 10$^{6}$ A m = 7.18 x 10$^{14}$ Dirac quanta, inherited from the monopole cloud and providing the handle for anti-Helmholtz confinement. The operating ratio g/$g_{\text{ext}}$ = 1.53 x 10$^{-10}$ places the black hole deeply sub-extremal; all Schwarzschild formulas apply.

Maldacena [17] showed magnetically charged black holes can exhibit enhanced Hawking emission via Landau-level quantization but this requires $m_W r_s$ >> 1. At the design point $m_W r_s$ = 0.605, so there is no enhancement.

\subsection{Why magnetic charge, not electric}
A micro black hole cannot sustain macroscopic electric charge. Even modest charge produces horizon potentials far exceeding $T_H$ and Schwinger pair production discharges the hole on nanosecond timescales. Magnetic charge is fundamentally different: discharging requires Schwinger production of monopole-antimonopole pairs, demanding B $\sim$ 10$^{49}$ T --- far beyond any physical environment. The magnetic charge is permanent. A pure magnetic field, regardless of strength, has B$^{2}$ - E$^{2}$ > 0 and E · B = 0; the vacuum is stable [18]. Only real photons undergo magnetic pair conversion, which is the mechanism by which direct Hawking photons convert to e$^{+}$e$^{-}$ pairs near the horizon.

\subsection{Particle census and loss budget}
Between the horizon and the SQM shell at R = 2 m, two transformations reshape the spectrum. All direct Hawking photons ($\sim$1\% of power) pair-convert within femtometers of the horizon. Quarks and gluons ($\sim$64\% of power) hadronize within $\sim$1 fm [19, 20]. After these processes, approximately 80\% of Hawking power resides in charged particles and $\sim$11\% in neutral hadrons. The total shell-intercepted fraction is $\sim$91\%. Irreducible losses are primary neutrinos (7.0\%), gravitons (0.1\%) and secondary neutrinos from stopped hadrons and muons ($\sim$10\%). Combined neutrino loss is $\sim$17\%, giving net captured power of 83\% --- 50,000 TW. The reference design: a 10$^{9}$ kg magnetically charged black hole radiating 60,200 TW, of which 50,000 TW is captured.

\section{THE SQM SHELL: A THERMAL WAVELENGTH CONVERTER}
The primary contribution of this paper is a CFL strange quark matter shell that absorbs the full GeV Hawking spectrum and re-emits it as magnetically directable keV electron-positron pairs, functioning as a thermal wavelength converter.

\subsection{The reflection problem}
Lee [2] showed that no known material reflects GeV-scale radiation. Abandoning reflection, he adopted titanium absorption at 33 km standoff, achieving a 46.5\% capture fraction and maximum velocity of 0.0001c. The deficiencies are compounding: the standoff distance is set by thermal survival, not any physics optimum and the absorbed energy re-radiates as omnidirectional blackbody infrared contributing nothing to thrust. The architecture below resolves both: a close-coupled thermal converter raises capture to 83\%, drops standoff from 33 km to 2 m and re-emits in a magnetically collimatable form.

\subsection{Architecture}
The shell is a spherical CFL enclosure with the black hole at its center. Shell parameters are listed in Table 3.

The shell is not a mirror. Its operating principle is thermal wavelength conversion: the inner surface absorbs the full GeV Hawking spectrum via nuclear-density stopping and the absorbed energy maintains a non-equilibrium electrosphere that re-emits power as keV electron-positron pairs. These pairs, being charged and sub-relativistic ($\beta$ = 0.332), are magnetically directable --- converting an omnidirectional GeV source into a collimated sub-MeV beam. The 2 m radius is set by the anti-Helmholtz coil geometry. The 3 pm thickness is set by the muon punch-through budget. The shell mass, $M_{\text{shell}}$ = $\rho$ · 4$\pi$R$^{2}$t = 6.0 x 10$^{7}$ kg, is approximately 3\% of total ship mass ($M_{\text{ship}}$ = 2$M_{\text{BH}}$ + $M_{\text{shell}}$ $\sim$ 2.06 x 10$^{9}$ kg).

The inner-surface electrosphere is the critical element. In normal operation, the shell's inner face is continuously bombarded by approximately 50,000 TW of GeV-scale radiation. This bombardment drives the electrosphere far out of thermal equilibrium, sustaining a pair-emission rate that exactly balances absorbed power at a self-consistent temperature $T_{\text{eq}}$ = 0.356 GK (kT = 30.7 keV). The re-emitted pairs carry away the absorbed energy at a wavelength roughly 10$^{5}$ times longer than the incident radiation --- hence "thermal wavelength converter."


\begin{table}[h]
\centering
\caption{SQM shell parameters.}
\label{tab:shell}
\small
\begin{tabular}{ll}
\hline
\textbf{Parameter} & \textbf{Value} \\
\hline
Inner radius $R$ & 2\,m \\
Thickness $t$ & 3\,pm (3000\,fm) \\
Density $\rho$ & $4 \times 10^{17}$\,kg/m$^3$ ($\sim$2$\times$ nuclear) \\
Mass $M_\text{shell}$ & $6.0 \times 10^7$\,kg (60,000\,t) \\
Surface area $A$ & $4\pi R^2 = 50.3$\,m$^2$ \\
Phase & Colour-flavour locked (CFL) \\
Inner surface state & Bombardment-maintained electrosphere \\
Outer surface state & Ground-state electrosphere \\
\hline
\end{tabular}
\end{table}

\subsection{Why CFL is not a superconducting mirror}
Manuel and Rajagopal [21] showed that in CFL, the ordinary photon mixes with the eighth gluon to form a rotated massless eigenstate that propagates freely: 99.75\% of the photon amplitude passes through unattenuated. CFL is a transparent insulator, not a superconductor. The thermal-converter architecture is the physically consistent alternative.

\subsection{Absorption: nuclear-density Bethe-Bloch stopping}
At $\rho$ = 4 x 10$^{17}$ kg/m$^{3}$, the mass thickness per femtometer is 40 g/cm$^{2}$. Applying the minimum-ionizing loss rate of 2 MeV/(g/cm$^{2}$) gives dE/dx = 80 MeV/fm. A caveat: the standard Bethe-Bloch formula is derived for atomic matter; at nuclear density the target is a degenerate quark plasma. Jet-quenching measurements at RHIC and LHC indicate loss rates 10---60 times higher in QGP than in cold nuclear matter [20], so 80 MeV/fm is a conservative lower bound.

At thickness t = 3000 fm, the total stopping budget is:

\begin{equation}
E_c = 80 \text{ MeV/fm} \times 3000 \text{ fm} = 240 \text{ GeV}
\end{equation}
All charged hadrons are stopped within the first few femtometers via the strong interaction. Neutral hadrons are captured within $\sim$1 fm of the inner surface. Electrons synchrotron-cool in the 227 T surface field and are stopped within $\sim$12 fm. Detailed species-by-species stopping analysis is available from the corresponding author on request.

Muon punch-through. Muons are the sole punch-through risk: their synchrotron cooling is suppressed relative to electrons by ($m_e$/$m_{\mu}$)$^{2}$ $\sim$ 2.4 x 10$^{-5}$ and they do not interact hadronically. Only Bethe-Bloch ionization stops them. Taus, despite similar properties, decay in flight: at $\gamma$ = 100 ($E_{\tau}$ $\sim$ 178 GeV), the laboratory decay length is 8.7 mm, well within the 2 m cavity, so taus deposit energy into secondary muons, hadrons and neutrinos before reaching the shell.

The muon species fraction is $f_{\mu}$/f = 4.67\% of total Hawking power. Including tau-decay secondary muons ($\sim$one-third of tau energy by branching ratio), the combined muon channel carries $\sim$9.3\% of $P_H$ ($\sim$5,600 TW). The Wien-tail escape fraction above $E_c$ = 240 GeV is:

\begin{equation}
\varepsilon(x_c) = \frac{e^{-x_c}(x_c^2 + 4x_c + 6)}{7\pi^4/120}, \quad x_c = E_c/kT
\end{equation}
At $x_c$ = 22.7: $\varepsilon$ = 1.5 x 10$^{-8}$. Shell thickness versus muon punch-through is detailed in Table 4.

The adopted 3,000 fm thickness yields only 84 MW of muon punch-through --- a negligible thermal loss (1.4 x 10$^{-9}$ of the total Hawking power). This energy is not deposited in the ship structure: the muons free-stream through the shell without radiating and decay in space hundreds of kilometers away. The 3,000 fm choice purchases six orders of magnitude of additional leakage suppression over the 640 fm alternative at a modest 4.7-fold increase in shell mass and negligible (<1.5\%) decrease in ship acceleration.


\begin{table}[h]
\centering
\caption{Shell thickness versus muon punch-through.}
\label{tab:punchthrough}
\small
\begin{tabular}{llllll}
\hline
\textbf{$t$ (fm)} & \textbf{$E_c$ (GeV)} & \textbf{$x_c$} & \textbf{$\varepsilon(x_c)$} & \textbf{$P_\text{punch}$ (muon)} & \textbf{$M_\text{shell}$ (t)} \\
\hline
640 & 51 & 4.84 & $6.8 \times 10^{-2}$ & 380\,TW & 12,900 \\
2,500 & 200 & 18.9 & $4.7 \times 10^{-7}$ & 2.6\,GW & 50,000 \\
3,000* & 240 & 22.7 & $1.5 \times 10^{-8}$ & 84\,MW & 60,000 \\
4,000 & 320 & 30.3 & $1.3 \times 10^{-11}$ & 74\,kW & 80,000 \\
5,000 & 400 & 37.8 & $1.0 \times 10^{-14}$ & 58\,W & 100,000 \\
\hline
\multicolumn{6}{l}{\footnotesize * Adopted design point.} \\
\end{tabular}
\end{table}

\subsection{Pair re-emission: the Usov mechanism}
Energy deposited in the inner $\sim$30 fm creates pairs at energies far exceeding the CFL gap $\Delta$ $\sim$ 100 MeV, bypassing gap-decoupling. The pair-emission flux follows the Usov [10, 22] formula:

\begin{equation}
f_{\pm} = 10^{39.2} \, T_9^3 \, \exp(-11.9/T_9) \, J(\zeta) \quad [\text{pairs cm}^{-2} \text{ s}^{-1}]
\end{equation}
where $T_9$ = T/(10$^{9}$ K), $\zeta$ = 20/$T_9$ and J($\zeta$) is the strong-field/thermal-flux interpolation function [10]. The formula was derived for neutron star surfaces at T > 10$^{9}$ K; the operating point sits at $T_{\text{eq}}$ = 3.56 x 10$^{8}$ K, slightly below the nominal domain but J($\zeta$) interpolates smoothly and remains well-defined at the operating point ($\zeta_{\text{eq}}$ = 56.2, J = 36.1).

A caveat regarding domain of applicability: the Usov formula was derived for neutron star surfaces at T > 10$^{9}$ K. Prakapenia and Vereshchagin [23] derived a $\sim$100-fold enhancement at T > 10$^{9}$ K using an updated kinetic treatment but their correction applies above the formula's nominal domain, not below it.

The self-consistent equilibrium is determined by balancing pair luminosity against absorbed bombardment power, divided by the magnetic-bottle capture fraction $f_{\text{cap}}$ = 12.1\%:

\[
L_{\text{pair}} = P_{\text{absorbed}}/(A \cdot f_{\text{cap}})
\]
\[
= 4.996 x 10^{16}/(50.3 x 0.121)
\]
\begin{equation}
L_\text{pair} = \frac{P_\text{abs}}{A \cdot f_\text{cap}} = 8.21 \times 10^{15} \text{ W/m}^2
\end{equation}
Setting $L_{\text{pair}}$ = 2$f_{±}$($m_e$ c$^{2}$ + kT) with self-consistent J($\zeta$) and solving yields:

\begin{equation}
T_\text{eq} = 0.356 \text{ GK}, \quad kT = 30.7 \text{ keV}, \quad \gamma = 1.060, \quad \beta = 0.332
\end{equation}
The robustness of this result deserves emphasis. Over a 100-fold variation in J (from 3.6 to 361), $\beta$ varies by only +/-3\% (from 0.323 to 0.342). The exponential sensitivity of the Usov rate to temperature acts as a thermostat: any perturbation in input power produces only a logarithmic shift in $T_{\text{eq}}$, which propagates as a percent-level change in $\beta$. The pair-exhaust velocity $\beta$ = 0.332 is robust to the residual Usov-formula uncertainty at the 3---5\% level.

\subsection{The three-layer outer-surface firewall}
The shell must absorb 50,000 TW on its inner surface and emit less than 1 MW from its outer surface --- a suppression factor of 5 x 10$^{10}$. Three mechanisms in series provide this:

(a) CFL gap suppression. Quark-quasiparticle excitation across the gap requires energy $\geq$ 2$\Delta$ $\sim$ 200 MeV. At kT = 30.7 keV, the Boltzmann factor is exp(-6510) $\sim$ 10$^{-2828}$.

(b) Massive-Goldstone suppression. CFL pseudo-Goldstone bosons ($\pi_{\text{CFL}}$, $K_{\text{CFL}}$) have masses m $\sim$ 5 MeV [24]. At kT = 30.7 keV: exp(-163) $\sim$ 10$^{-71}$.

(c) Massless-Goldstone anomaly emission. The CFL U(1)\_B Goldstone ($\eta$'-like mode) couples to photons via the Wess-Zumino anomaly. The leading process is $\eta$' + $\eta$' $\to$ 2$\gamma$ via t- and u-channel exchange, carrying two anomaly insertions. The squared amplitude scales as |M|$^{2}$ $\sim$ $\alpha ^{4}$T$^{4}$/$F_{\pi} ^{4}$ (not $\alpha ^{2}$, because two anomaly vertices are required). The bulk emissivity is [24]:

\begin{equation}
P/V \approx \frac{\alpha^4 T^9}{\pi^4 F_\pi^4}
\end{equation}
where $F_{\pi}$\_CFL $\sim$ 100 MeV. At $T_{\text{bulk}}$ = 30.7 keV, this yields P/V = 2.2 x 10$^{14}$ W/m$^{3}$. Integrated over shell volume V = A · t = 1.5 x 10$^{-10}$ m$^{3}$, the raw bulk emission is 34 kW. These anomaly-produced photons have characteristic energy $\sim$2kT $\sim$ 60 keV --- well below the outer-electrosphere plasma frequency. The outer electrosphere is a degenerate ultrarelativistic electron gas whose plasma frequency at conservative surface chemical potential $\mu_e$ = 9 MeV [25, 26] is:

\begin{equation}
\omega_p = \sqrt{\frac{4\alpha}{3\pi}} \, \mu_e = 0.50 \text{ MeV}
\end{equation}
The Wien-tail escape fraction above $\omega_p$ ($x_c$ = $\omega_p$/kT = 16.3) is $\varepsilon$ = 5 x 10$^{-6}$. Net leakage: 34 kW x 5 x 10$^{-6}$ = 0.17 W. Even with 10$^{3}$ uncertainty in the anomaly prefactor, leakage remains below 170 W --- a factor of 6 x 10$^{3}$ below the 1 MW off-axis budget.

The combination of three layers is the load-bearing firewall: (a) kills quark-quasiparticle channels by 10$^{-2828}$; (b) kills massive-Goldstone channels by 10$^{-71}$; (c) produces the only surviving emission channel at 34 kW, then the outer electrosphere attenuates the escaping fraction by 5 x 10$^{-6}$. No single layer is asked to do the entire job.

\subsection{Thermal leakage summary}
The shell converts 50,000 TW of absorbed GeV radiation to keV pairs with total electromagnetic leakage below 1 kW at the design point. The only non-exhaust energy losses are muon punch-through (84 MW, free-streaming away from the ship), neutrinos ($\sim$10,200 TW, non-interacting) and gravitons ($\sim$60 TW). No radiator or active thermal rejection system is required --- all captured energy exits through the exhaust bore by design.

\subsection{Coulomb barrier and cosmological safety}
The Coulomb barrier between incident nuclei and the SQM surface is 1---5 MeV depending on Z. At room temperature (kT = 0.025 eV), tunneling probability is $\sim$10$^{-300}$. SQM is effectively inert below $\sim$10$^{10}$ K. The ship structure contacts only the outer electrosphere --- a layer of ordinary electrons --- not the bare quark surface.

\section{CONFINEMENT, CAPTURE AND EXHAUST}
A single anti-Helmholtz coil pair solves two problems simultaneously: it confines the magnetically charged black hole near the shell's geometric center and creates the magnetic bottle that captures re-emitted pairs and funnels them into the exhaust nozzle.

\subsection{Anti-Helmholtz geometry}
The system consists of two superconducting coils of radius $R_c$ = 2 m at axial positions z = +/-d = +/-2 m, carrying opposing currents I = 1.36 x 10$^{9}$ A. The on-axis magnetic field is:

\begin{equation}
B_z(z) = \frac{\mu_0 I R_c^2}{2} \left[ \frac{1}{(R_c^2 + (d-z)^2)^{3/2}} - \frac{1}{(R_c^2 + (d+z)^2)^{3/2}} \right]
\end{equation}
which vanishes at the center (z = 0) by symmetry and produces a saddle-point topology: the field magnitude increases in both axial and radial directions away from the center. The axial gradient at center is:

\begin{equation}
\left.\frac{dB_z}{dz}\right|_{z=0} = \frac{3\mu_0 I R_c^2 d}{(R_c^2 + d^2)^{5/2}} = 227 \text{ T/m}
\end{equation}
Field magnitudes at the shell surface (R = 2 m): 390 T at poles and 227 T at equator (from div B = 0 giving $B_r$ = -(r/2) $dB_z$/dz). These fields require far-future superconductors at sustained currents of 1.36 x 10$^{9}$ A with peak conductor fields of $\sim$130 T and hoop stresses of $\sim$6.7 GPa. Stored energy is approximately 50 GJ in a coil mass of $\sim$1,000 kg.

\subsection{Black hole confinement: equilibrium and stability}
The force on the black hole's magnetic charge g along the symmetry axis is $F_{\text{mag}}$ = g($dB_z$/dz)$\delta$ for small displacements $\delta$ from the field zero. In the accelerating reference frame of the ship, the black hole experiences an inertial pseudo-force $F_{\text{inertial}}$ directed aft. Steady-state equilibrium requires:

\begin{equation}
g \cdot \frac{dB}{dz} \cdot \delta_\text{eq} = F_\text{thrust} \cdot \frac{M_\text{BH}}{M_\text{ship}}
\end{equation}
With $F_{\text{thrust}}$ = 55.3 MN, $M_{\text{BH}}$/$M_{\text{ship}}$ = 10$^{9}$/(2.06 x 10$^{9}$) = 0.485 and $\delta_{\text{eq}}$ = 100 mm:

\[
g = F_{\text{inertial}}/((dB/dz) \cdot \delta{}_eq)
\]
\[
= 2.68 x 10^{7}/(227 x 0.100)
\]
\begin{equation}
g = 1.182 \times 10^6 \text{ A\,m}
\end{equation}
corresponding to $N_D$ = 7.18 x 10$^{14}$ Dirac quanta.

Axial stability is passive. A small displacement from $\delta_{\text{eq}}$ produces a restoring force directed toward the equilibrium point. The gradient is linear in z near the center and the pseudo-force is constant, so the dynamics are those of a simple harmonic oscillator --- inherently stable without feedback.

Lateral stability is not passive. By Earnshaw's theorem, an axially symmetric static magnetic field cannot provide simultaneously stable equilibrium in all three dimensions. The radial field component near the axis produces a destabilizing force with negative spring constant $k_r$ = -g|dB/dz|/2 = -1.3 x 10$^{8}$ N/m. The natural instability timescale is:

\begin{equation}
\tau_\text{lat} = \sqrt{\frac{M_\text{BH}}{|k_r|}} \approx 2.8 \text{ s}
\end{equation}
Active PD control with $\sim$1 Hz bandwidth (ten times the instability rate) stabilizes the system. Position readout exploits the black hole's own monopole field: at 2 m the field is 30 T with gradient $\sim$30 T/m, giving $\sim$30 pm precision at a 1 nT magnetometer noise floor. The actuator is realized by asymmetric modulation of the coil currents ($\sim$1\% of baseline). This is standard control engineering, not a fundamental physics challenge.

\subsection{Magnetic bottle pair capture}
Pairs emerge from the inner electrosphere with $\gamma$ = 1.060, $\beta$ = 0.332. Their Larmor radius in the 227 T equatorial field is:

\[
r_L = \gamma{}\beta{}m_e c/(eB)
\]
\[
= 1.060 x 0.332 x 9.11 x 10^{-31} x 3 x 10^{8}/(1.602 x 10^{-19} x 227)
\]
\begin{equation}
r_L = \frac{\gamma \beta m_e c}{eB} = 2.65 \text{ \textmu m}
\end{equation}
which is 8 x 10$^{-7}$ of the shell radius. The pairs are deeply magnetized everywhere within the cavity and follow field lines with negligible cross-field drift.

The anti-Helmholtz geometry creates a magnetic bottle. The field strength increases from 227 T at the equator to 1,000 T at the exhaust bore throat. A pair launched with pitch angle $\alpha$ relative to the local field is reflected at the mirror point where sin$^{2}$($\alpha$) = $B_{\text{loc}}$/$B_{\text{throat}}$. The escape solid-angle fraction for isotropic equatorial emission is:

\[
f_{\text{cap}} = 1 - cos(arcsin sqrt(B_{\text{loc}}/B_{\text{throat}}))
\]
\begin{equation}
f_\text{cap} = 1 - \cos\!\left(\arcsin\sqrt{\frac{B_\text{loc}}{B_\text{throat}}}\right) = 12.1\%
\end{equation}
The energy balance uses this conservative equatorial value, treating all emission as equatorial. At the poles, where $B_{\text{loc}}$ = 390 T, the capture fraction rises to 21.9\%.

Pairs that miss the loss cone on first emission are magnetically reflected and bounce between mirror points with cavity transit time $\tau_b$ = R/v $\sim$ 2/(0.332c) = 19 ns. In a collisionless bottle these pairs would be permanently trapped. Coulomb scattering between bouncing pairs randomizes the pitch angle: the Mott cross-section for e$^{+}$/e$^{-}$ scattering at 30.7 keV is $\sigma_M$ $\sim$ 5.8 x 10$^{-27}$ m$^{2}$ and the steady-state cavity density builds to n* $\sim$ 10$^{24}$ m$^{-3}$, at which point collisional pitch-angle diffusion refills the loss cone on a timescale of approximately six bounces ($\sim$120 ns). Every emitted pair eventually finds the loss cone and escapes through the bore. The 12.1\% first-pass figure used in the energy balance coincides with the steady-state escape rate per emission event in this collisional regime. Synchrotron losses during multi-bounce residence are negligible: $\tau_{\text{sync}}$ = 1.28 ms at $\gamma$ = 1.060 in 227 T, exceeding the $\sim$120 ns residence time by a factor of 10$^{4}$.

\subsection{Exhaust nozzle}
The exhaust bore is a 0.3 mm diameter aperture at the aft pole, where the field reaches 1,000 T. Pairs enter the bore with Larmor radius 0.63 $\mu$m and proceed into an 820 m exponentially tapered magnetic nozzle (1,000 T throat to 30 $\mu$T exit). The adiabatic invariant $\mu$ = $p_{\text{perp}} ^{2}$/(2mB) is conserved throughout. Energy conservation with constant $\mu$ gives:

\[
sin^{2}(\theta{}_exit) = B_{\text{exit}}/B_{\text{throat}} = 3 x 10^{-5}/10^{3} = 3 x 10^{-8}
\]
\begin{equation}
\sin^2\theta_\text{exit} = \frac{B_\text{exit}}{B_\text{throat}} = 3 \times 10^{-8}, \quad \theta_\text{exit} \approx 0.010°
\end{equation}
Thrust-divergence loss is (1 - cos $\theta_{\text{exit}}$) $\sim$ 10$^{-8}$ --- entirely negligible. The exhaust beam is collimated to within a hundredth of a degree. The forward refueling port is SQM-capped during cruise, reflecting all pairs aft; during refueling, geometric leakage through the open port (bore solid angle = 1.56 x 10$^{-9}$ of 4$\pi$) is less than 100 kW.

\section{FORMATION AND MONOPOLE PRODUCTION}
\subsection{BPS monopole gravitational collapse}
The black hole is formed by gravitational collapse of a cloud of same-sign BPS 't Hooft-Polyakov monopoles. Each monopole carries one Dirac quantum $g_D$ = 1.646 x 10$^{-9}$ A m with BPS mass $m_M$ = 1.77 x 10$^{17}$ GeV/c$^{2}$ = 3.15 x 10$^{-10}$ kg [27, 28, 8, 9]. The required $N_D$ = 7.18 x 10$^{14}$ monopoles yield cloud mass $M_{\text{cloud}}$ = 226 tonnes.

The critical physics enabling gravitational collapse is the Bogomolny-Prasad-Sommerfield limit. In a gauge theory with Higgs self-coupling $\lambda$, same-sign monopoles experience two competing long-range forces: repulsive 1/r$^{2}$ magnetic Coulomb and attractive scalar Yukawa exchange. In general, magnetic repulsion dominates by a factor of order 10$^{5}$. In the BPS limit ($\lambda$ $\to$ 0), Manton [29] and Gibbons and Manton [30] proved that the scalar attraction exactly cancels magnetic repulsion at all separations. This cancellation is not fine-tuning: in GUTs with extended N = 2 SUSY in the monopole sector, $\lambda$ = 0 is protected to all loop orders by non-renormalization theorems.

With electromagnetic forces canceled, the cloud is pressureless dust and collapse follows Oppenheimer-Snyder dynamics. For a uniform-density cloud released from rest at $R_{\text{cloud}}$ = 10 mm (inter-monopole spacing $\sim$2.5 $\mu$m, density $\rho$ = 5.4 x 10$^{10}$ kg/m$^{3}$), the free-fall time is:

\begin{equation}
t_\text{ff} = \sqrt{\frac{3\pi}{32 G\rho}} = 0.29 \text{ s}
\end{equation}
The resulting black hole inherits the full magnetic charge g = $N_D$ x $g_D$ = 1.182 x 10$^{6}$ A m --- precisely the operating charge for anti-Helmholtz confinement.

The BPS force cancellation imposes a constraint on the Higgs self-coupling: the scalar Higgs force has range $r_H$ = $\hbar$/($m_H$ c) and for $r_H$ to exceed the inter-monopole spacing at initial cloud density ($\sim$2.5 $\mu$m), the condition requires $\lambda$ < 10$^{-53}$ at the GUT scale. This is extraordinary in isolation but N = 2 SUSY protection enforces $\lambda$ = 0 exactly, satisfying it trivially.

The same-sign requirement is critical: monopole-antimonopole pairs experience doubly attractive forces and would annihilate prematurely, releasing 2$m_M$ c$^{2}$ $\sim$ 5.7 x 10$^{7}$ J of gauge radiation per pair. Prior to collapse, all 7.18 x 10$^{14}$ monopoles are stored individually in niobium Meissner traps --- hollow superconducting spheres of inner radius $\sim$100 nm and mass 43 fg per trap (total trap mass 0.5 kg). In interstellar space the ambient 2.7 K background is well below niobium's $T_c$ = 9.3 K and traps remain superconducting indefinitely without active refrigeration. Release is achieved by heating above $T_c$ via microwave pulse.

Seed black hole properties. The 226-tonne black hole has kT $\sim$ 46.8 TeV, P $\sim$ 1.4 x 10$^{12}$ TW and unfed lifetime of 5.3 milliseconds. Growing this seed to operating mass is the bootstrap task.

\subsection{Core-overlap exponential breeding}
A single seed pair is multiplied to 7.18 x 10$^{14}$ monopoles via exponential breeding in core-overlap collisions.

The dual coupling. By the Dirac quantization condition, the minimum magnetic charge is $g_D$ = 2$\pi$/e, giving magnetic fine-structure constant $\alpha_m$ = $g_D ^{2}$/(4$\pi$) = 1/(4$\alpha_e$) --- the Olive-Witten [12] duality relation. At low energies ($\alpha_e$ = 1/137), $\alpha_m$ $\sim$ 34, deeply non-perturbative. Inside a monopole core, the Higgs VEV vanishes and GUT symmetry is locally restored. The relevant coupling becomes $\alpha_m$ = 1/(4$\alpha_{\text{GUT}}$) $\sim$ 6.25. The suppression for pair production is exp(-2$\pi$/$\alpha_m$) $\sim$ 0.37, essentially unsuppressed when sufficient energy is available.

The breeding reaction. When two same-sign monopoles collide with impact parameter b < $r_c$ $\sim$ 2.8 x 10$^{-32}$ m (the core radius, set by $\hbar$c/$M_X$), their cores overlap and GUT symmetry is restored. The reaction M + M $\to$ M + M + M + $M_{\text{bar}}$ proceeds with O(1) probability, conserving topological charge (+2$g_D$ $\to$ +3$g_D$ - $g_D$ = +2$g_D$). The energy threshold is 4$m_M$ c$^{2}$; at collision Lorentz factor $\gamma$ = 3 per beam, $E_{\text{CM}}$ = 6$m_M$ c$^{2}$ provides a kinetic surplus of 2$m_M$ c$^{2}$ above threshold.

Generation count. Each generation takes N monopoles, forms N/2 collision pairs and produces N/2 additional monopoles (plus N/2 antimonopoles stored separately). Population grows by factor 1.5 per generation. Starting from one pair, reaching 7.18 x 10$^{14}$ requires lo$g_{1.5}$(2 x 7.18 x 10$^{14}$) = 86 generations. The breeding takes place in SQM-clad synchrotron rings at $\sim$7 x 10$^{7}$ T bending field (circumference $\sim$14 m for $\gamma$ = 3 monopoles). With realistic beam-loading constraints, the total breeding time ranges from hours to years.

\subsection{Seed production: SQM electron linac}
Producing the first monopole-antimonopole pair is the most speculative element. Drukier and Nussinov [31] showed that production of 't Hooft-Polyakov monopoles in collisions of pointlike Standard Model particles is exponentially suppressed:

\begin{equation}
\Gamma \propto \exp(-c/\alpha_\text{GUT}), \quad c \approx 4
\end{equation}
At $\alpha_{\text{GUT}}$ = 0.04, this gives exp(-100) $\sim$ 10$^{-43}$ per collision at threshold. The suppression arises from the topological mismatch between an incoming pointlike wavefunction and the smooth classical soliton profile. This result is not in dispute at threshold.

The seed-production apparatus is an SQM electron linac: two opposing 2,500 km CFL-bore tubes at an accelerating gradient of 10$^{21}$ V/m (6---50x below the CFL critical field), each beam reaching 2.5 x 10$^{18}$ GeV. Head-on collisions produce $E_{\text{CM}}$ = 5 x 10$^{18}$ GeV --- ten times the pair threshold $E_{\text{thr}}$ = 2$m_M$ = 3.5 x 10$^{18}$ GeV. At 10$^{6}$ Hz repetition and conservative 10$^{-7}$ production probability, first pair arrives in $\sim$10 seconds.

Four arguments support the claim that Drukier-Nussinov suppression collapses at 10x threshold:

(1) Classical GUT symmetry restoration. At 10x threshold, the collision deposits $E_{\text{CM}}$ = 3.5 x 10$^{19}$ GeV in a volume of order $r_c ^{3}$, producing energy density $\sim$ 1.3 x 10$^{66}$ GeV$^{4}$. The GUT Higgs potential barrier has height $\lambda$v$^{4}$/4; for any $\lambda$ < 505 (including all physically relevant values and in particular the BPS limit $\lambda$ = 0), the collision energy density exceeds the barrier. The GUT vacuum is classically restored, monopoles become perturbative excitations and the topological mismatch underlying Drukier-Nussinov does not apply.

(2) Holy-grail function F < 0. The energy-dependent suppression can be written as exp(-4$\pi$/$\alpha_{\text{GUT}}$ x F($\varepsilon$)), where $\varepsilon$ = $E_{\text{CM}}$/$E_{\text{thr}}$ - 1 and F($\varepsilon$) $\sim$ 1 - (9/8)$\varepsilon ^{4}$ᐟ$^{3}$ [32]. At threshold ($\varepsilon$ = 0), F = 1 and the full Drukier-Nussinov suppression applies. F crosses zero at $\varepsilon$ $\sim$ 0.9 ($E_{\text{CM}}$ $\sim$ 1.9$E_{\text{thr}}$). At the linac's operating point ($\varepsilon$ = 9), F = -20 --- unphysical, indicating the system is deep in the non-perturbative over-barrier regime where tunneling no longer applies.

(3) Kibble freeze-out. Once the collision classically restores GUT symmetry in a region of radius $\geq$ $r_c$, monopoles form by the Kibble mechanism [33] as the bubble cools and symmetry breaks again. The number of defects is N $\sim$ $V_{\text{bubble}}$/$\xi ^{3}$ $\sim$ 1, where $\xi$ $\sim$ $r_c$ is the correlation length. The production cross-section is geometric:

\begin{equation}
\sigma \sim \pi r_c^2 \sim 2.4 \times 10^{-63} \text{ m}^2
\end{equation}
exceeding the pointlike cross-section by $\sim$10$^{5}$ because $r_c$ far exceeds the de Broglie wavelength. No exponential suppression.

(4) BPS-specific: no barrier at $\lambda$ = 0. In the BPS limit, the Higgs potential V($\varphi$) vanishes identically along D-flat directions. The field space is a smooth manifold (the monopole moduli space) with no energy barriers between the perturbative vacuum and the soliton sector. The Drukier-Nussinov suppression exp(-c/$\alpha$) arises from the WKB integral through this barrier; when the barrier height is zero, the integral vanishes and the suppression factor is exp(0) = 1 at any energy. The only costs are kinematic (providing rest mass 2$m_M$) and gradient (constructing the soliton profile). The gradient energy $E_{\text{gradient}}$ $\sim$ $v_{\text{GUT}}$/$g_{\text{YM}}$ $\sim$ 1.4 x 10$^{16}$ GeV is exceeded by the collision energy by a factor of 250.

The combined production probability per linac collision is estimated at 10$^{-7}$ (conservative) to O(1) (Kibble freeze-out per energy-depositing collision). The factor of 10$^{-7}$ likely reflects the kinematic bottleneck --- the fraction of e$^{-}$e$^{-}$ collisions that concentrate sufficient energy in a core-sized volume --- rather than any residual tunneling suppression. No lattice calculation exists for monopole production at finite collision energy in any specific GUT and non-perturbative corrections to the holy-grail function remain uncomputed. Both aggressive and conservative estimates are presented; even the conservative case gives $\sim$10 seconds to first pair at 10$^{6}$ Hz.

If the linac proves impractical, the fallback is passive capture of a primordial GUT monopole from the cosmological Kibble-mechanism residue, using a $\sim$10 km$^{2}$ superconducting flux-concentrator sail over $\sim$10$^{4}$ years. The current observational upper bound (MACRO 2002: $\varphi$2 < 1.4 x 10$^{-16}$ cm$^{-2}$ s$^{-1}$ sr$^{-1}$) [37] permits this but does not guarantee it. The downstream physics is entirely independent of how the first pair is obtained.

\subsection{Bootstrap: seed to operating mass}
The 226-tonne seed has an unfed lifetime of 5.3 milliseconds. It must be grown to 10$^{9}$ kg by feeding matter faster than it evaporates. The bootstrap proceeds by continuous railgun injection of pellets at $\beta$ = 0.33. Each pellet contributes both rest mass and kinetic energy; the kinetic fraction is ($\gamma$ - 1)/$\gamma$ $\sim$ 0.061. A temporary recovery shell captures a fraction $f_{\text{cap}}$ of Hawking radiation and converts it to pellet kinetic energy. The net mass-growth rate is:

\begin{equation}
\frac{dM}{dt} = A \cdot \frac{P(M)}{c^2}
\end{equation}
\begin{equation}
A = \frac{\gamma f_\text{cap}}{(\gamma - 1)/\gamma} - 1
\end{equation}
where P(M) is the Hawking power at mass M. For $f_{\text{cap}}$ = 0.5, the amplification factor is A $\sim$ 8.5 --- the hole gains 8.5 kg per kilogram lost. The growth equation M$^{2}$ dM = (A$\hbar$c$^{4}$f/G$^{2}$) dt integrates to give:

\begin{equation}
t_\text{boot} = \frac{G^2(M_f^3 - M_0^3)}{3A\hbar c^4 f}
\end{equation}
Numerical integration with the full f(M) tabulation, accounting for Standard Model threshold crossings as kT decreases from 46.8 TeV to 10.6 GeV, yields approximately 2 years (1.6 yr analytic; 1.85 yr numerical with $f_{\text{cap}}$ = 0.5). The bootstrap is dominated by the late phase: M grows from 226 tonnes to 10$^{8}$ kg in minutes but the final octave (5 x 10$^{8}$ to 10$^{9}$ kg) takes over a year because dM/dt scales as 1/M$^{2}$. The system is mass-positive provided $f_{\text{cap}}$ > 0.061. At $f_{\text{cap}}$ = 0.3, $t_{\text{bootstrap}}$ $\sim$ 2.9 yr; at $f_{\text{cap}}$ = 0.7, $t_{\text{bootstrap}}$ $\sim$ 1.1 yr.

\section{PERFORMANCE}
The three physics elements --- the CFL SQM thermal wavelength converter, the anti-Helmholtz confinement and exhaust system and the monopole-collapse formation mechanism --- combine into a propulsion architecture with three properties no other concept offers simultaneously.

\subsection{Thrust is continuously adjustable}
At the reference design point (M = 10$^{9}$ kg), thrust is:

\begin{equation}
F = \frac{\beta P_\text{cap}}{c} = 55.3 \text{ MN}
\end{equation}
with specific impulse $I_{\text{sp}}$ = $\beta$c/g₀ = 1.01 x 10$^{7}$ s and photon-rocket efficiency Fc/$P_H$ = 27.5\%.

The black hole mass is a free parameter. Since $P_H$ scales as 1/M$^{2}$, the same architecture at lower mass produces quadratically more thrust. Shell thickness scales linearly with kT (to maintain the stopping budget $E_c$ = 80 MeV/fm x $t_{\text{shell}}$), increasing shell mass; coils, bore and nozzle geometry are unchanged. Thrust scaling with black hole mass is detailed in Table 5.

All profiles carry fuel from launch. Mass ratios and peak velocities are computed from the relativistic Tsiolkovsky equation:

\begin{equation}
v_\text{peak} = c \tanh\!\left(\frac{v_\text{eff}}{c} \ln R_\text{leg}\right)
\end{equation}
with total mass ratio $R_{\text{total}}$ = $R_{\text{leg}} ^{2}$ (symmetric boost/decelerate) and effective exhaust velocity $v_{\text{eff}}$ = 0.83 x $\beta_e$ x c = 0.276c. The factor 0.83 accounts for the 17\% neutrino loss: fuel feeds the black hole at rate $P_H$/c$^{2}$ but only 83\% of the radiated mass contributes to thrust.

There is no physics ceiling on thrust. The practical limits are: (a) shell thickness grows linearly with kT, increasing shell mass quadratically as M decreases; (b) fuel mass grows exponentially with $v_{\text{peak}}$ via the rocket equation; (c) fuel rate scales as 1/M$^{2}$, demanding larger supply infrastructure. These are engineering constraints, not physics constraints. The architecture is identical at every point in the table --- the same shell material, the same confinement geometry, the same exhaust mechanism. The black hole mass is the single design dial.


\begin{table}[h]
\centering
\caption{Thrust scaling with black hole mass.}
\label{tab:thrust}
\small
\begin{tabular}{llllll}
\hline
\textbf{$M$ (kg)} & \textbf{$kT$ (GeV)} & \textbf{$F$ (MN)} & \textbf{Shell} & \textbf{Fuel rate} & \textbf{Alpha Cen} \\
\hline
$10^9$* & 10.6 & 55 & 3\,pm, 60\,kt & 0.67\,kg/s & 79\,yr, $v_\text{peak}$\,0.07$c$, $R{=}1.7$ \\
$2 \times 10^8$ & 53 & 1,375 & 10\,pm, 200\,kt & 17\,kg/s & 19\,yr, $v_\text{peak}$\,0.30$c$, $R{=}9.4$ \\
$5 \times 10^7$ & 212 & 21,800 & 30\,pm, 600\,kt & 264\,kg/s & 11\,yr, $v_\text{peak}$\,0.45$c$, $R{=}34$ \\
\hline
\multicolumn{6}{l}{\footnotesize * Reference design.} \\
\end{tabular}
\end{table}

\subsection{The fuel is any matter}
The black hole converts mass to Hawking radiation at mc$^{2}$ per kilogram --- the theoretical maximum for any energy source. The Hawking spectrum depends only on M and g, not on infalling composition. Hydrogen, carbon, iron, rock, ice, regolith, water, organic waste --- all produce identical thrust.

At the reference design, fueling is 670 g/s (21,140 tonnes/year --- roughly two small asteroids per century, of any composition). At high-thrust (M = 2 x 10$^{8}$ kg), 17 kg/s (540,000 tonnes/year, equivalent to one 500-meter asteroid per year). No other propulsion concept has this property: antimatter requires antimatter, fusion requires hydrogen isotopes, fission requires actinides. The black hole drive requires mass --- the most abundant resource in the universe.

\subsection{Thermal management is solved}
In every other high-thrust concept --- nuclear thermal, antimatter, fusion --- thermal management is the scaling limit. More thrust means more waste heat; the engine melts before mission parameters are met. The SQM shell sidesteps this entirely. The thermal budget has three distinct zones:

Inner surface (bombardment face). The inner electrosphere absorbs 50,000 TW over 50 m$^{2}$ --- an intensity of 10$^{15}$ W/m$^{2}$. This is by design: the energy deposits in the inner $\sim$30 fm and is immediately re-emitted as Usov pairs. The Usov rate rises so steeply with temperature that even a 100-fold increase in input power (e.g. at M = 10$^{8}$ kg) raises $T_{\text{eq}}$ only logarithmically to $\sim$35 keV. The inner surface is a thermostat, not a heat sink.

CFL bulk. The CFL gap ($\Delta$ $\sim$ 100 MeV) blocks thermal transport from the inner skin to the bulk [13]. At kT = 30.7 keV --- a factor of 3,000 below the gap --- quark-quasiparticle conduction is suppressed by exp(-2$\Delta$/kT) = 10$^{-2828}$. The bulk stays cold.

Outer surface (three-layer firewall). Three mechanisms suppress outward emission by a combined factor of $\sim$10$^{10}$. Net electromagnetic leakage is below 1 kW at any operating point.

All absorbed energy exits through the aft bore as pair exhaust. There is no radiator, no cooling loop, no thermal rejection system. The thrust dial is real because none of the three thermal zones limits with increasing power. The engine cannot overheat.

\subsection{Mission profiles}
All profiles assume fuel carried from launch (any matter), symmetric boost/decelerate trajectories and the relativistic Tsiolkovsky equation with $v_{\text{eff}}$ = 0.276c. Mission profiles are detailed in Table 6.

A mass ratio of 10 means carrying 9x the dry mass in any solid material --- roughly one 1-km asteroid per 10$^{8}$ kg of dry mass. The Sprint design to Alpha Centauri (11 years, $v_{\text{peak}}$ = 0.45c) requires 33x dry mass in fuel --- at $M_{\text{dry}}$ $\sim$ 1.2 x 10$^{9}$ kg (BH + shell + cargo), the fuel mass is $\sim$4 x 10$^{10}$ kg, a single 300-meter rocky asteroid.

Unfed delta-v (burning the black hole without refueling). For a ship with $M_{\text{ship}}$ = 2$M_{\text{BH}}$ + $M_{\text{shell}}$, evaporating one $M_{\text{BH}}$ gives mass ratio $\sim$1.49. The relativistic Tsiolkovsky equation yields $\Delta$v = c tanh(0.276 x ln 1.49) $\sim$ 0.11c. In practice this is limited by the shell's stopping budget: as the hole shrinks, kT rises past the 240 GeV stopping budget at M < 10$^{7}$ kg.


\begin{table}[h]
\centering
\caption{Mission profiles with relativistic Tsiolkovsky accounting.}
\label{tab:missions}
\small
\begin{tabular}{llllll}
\hline
\textbf{Target} & \textbf{Design} & \textbf{$v_\text{peak}$} & \textbf{Mass ratio $R$} & \textbf{Fuel ($\times M_\text{dry}$)} & \textbf{Trip time} \\
\hline
Alpha Cen (4.37\,ly) & Reference ($10^9$\,kg) & 0.07$c$ & 1.7 & 0.7$\times$ & 79\,yr \\
Alpha Cen & High-thrust ($2{\times}10^8$\,kg) & 0.30$c$ & 9.4 & 8.4$\times$ & 19\,yr \\
Alpha Cen & Sprint ($5{\times}10^7$\,kg) & 0.45$c$ & 34 & 33$\times$ & 11\,yr \\
10\,ly & Reference & 0.10$c$ & 2.0 & 1.0$\times$ & 113\,yr \\
10\,ly & High-thrust & 0.35$c$ & 14 & 13$\times$ & 36\,yr \\
100\,ly & High-thrust & 0.50$c$ & 54 & 53$\times$ & 220\,yr \\
\hline
\end{tabular}
\end{table}

\subsection{Electromagnetic quietness}
A convenient consequence of the shell's thermal architecture is that the drive produces virtually no off-axis electromagnetic emission. Conventional high-thrust propulsion concepts --- nuclear thermal, fusion, antimatter --- require large radiator arrays to reject waste heat, producing thermal signatures proportional to thrust. The SQM shell eliminates this: all captured energy exits through the exhaust bore, and the three-layer firewall limits total off-axis electromagnetic leakage to less than 1 kW. The residual non-electromagnetic losses (84 MW of free-streaming muons, $\sim$8,500 TW of neutrinos) are non-interacting and deposit no energy in the surrounding environment. The drive is, as an incidental consequence of its thermal design, among the quietest high-thrust propulsion concepts conceivable.

The SQM thermal-converter shell is the enabling innovation. Crane and Westmoreland [3] proposed the concept but left energy capture unspecified; Lee [2] demonstrated that conventional materials fail, achieving at most 0.0001c. The CFL shell raises capture from 46.5\% to 83\%, eliminates the 33 km standoff (the shell is 2 m from the black hole) and introduces self-regulating thermal management that removes the thermal scaling limit on thrust. A comparison to prior work is provided in Table 7.


\begin{table}[h]
\centering
\caption{Comparison to prior work.}
\label{tab:comparison}
\small
\begin{tabular}{lllll}
\hline
 & \textbf{Crane \& W. 2009} & \textbf{Lee 2015} & \textbf{Present (ref.)} & \textbf{Present (sprint)} \\
\hline
Capture mechanism & \parbox[t]{2.5cm}{``Electron gas mirror'' (unspecified)} & \parbox[t]{2.2cm}{Titanium absorption} & \parbox[t]{2.5cm}{CFL SQM thermal converter} & Same \\[6pt]
Capture fraction & Assumed $\sim$100\% & 46.5\% & 83\% & 83\% \\
Alpha Cen trip & Not calculated & Not achievable & 79\,yr at 0.07$c$ & 11\,yr at 0.45$c$ \\
Thrust & Not calculated & --- & 55\,MN & 21,800\,MN \\
$I_\text{sp}$ & Not calculated & --- & $1.01 \times 10^7$\,s & Same \\
Fuel & Not specified & N/A & Any matter & \parbox[t]{2cm}{Any matter ($33\times M_\text{dry}$)} \\[4pt]
Thermal limit & Not addressed & \parbox[t]{2.2cm}{Titanium melting (33\,km standoff)} & \parbox[t]{2.5cm}{None (Usov self-regulation)} & Same \\[6pt]
Confinement & ``Particle beams'' & None & Anti-Helmholtz & Same \\
Formation & \parbox[t]{2.5cm}{Gamma-ray kugelblitz} & N/A & \parbox[t]{2.5cm}{BPS monopole collapse} & Same \\
\hline
\end{tabular}
\end{table}

\section{DISCUSSION}
\subsection{The SQM shell as the enabling physics}
The central result is that a CFL shell transforms black hole propulsion from a theoretical curiosity into a quantitatively viable architecture. Lee [2] showed no conventional material can reflect the GeV Hawking spectrum. The SQM thermal wavelength converter eliminates both the reflection problem and the 33 km thermal-standoff constraint, raising capture from 46.5\% to 83\% and converting the 0.0001c ceiling into a 0.07c---0.45c performance range. Combined with continuously adjustable thrust (1/M$^{2}$ scaling) and universal fuel, the shell converts a single-point concept into an interstellar-capable propulsion architecture.

\subsection{Assumption sensitivity}
The design rests on six assumptions (Table 1) and the consequences of each failing independently merit explicit statement. If Hawking radiation does not exist (Assumption 1), the entire concept is void --- there is no power source. This assumption enjoys near-universal theoretical consensus [5] but remains observationally unconfirmed. If GUT monopoles do not exist (2) or the BPS limit does not hold (3), the formation mechanism fails; the SQM shell, confinement system and exhaust architecture survive intact and would function with any alternative micro black hole source. If the Bodmer-Witten hypothesis is false (4) or CFL is not the ground state at the operating chemical potential (5), there is no shell material and the design reverts to the 0.0001c bound with conventional absorption. If monopole breeding fails (6), the inventory must be obtained by other means --- primordial-flux capture over $\sim$10$^{4}$ years with a superconducting sail, which is not engineering-feasible but is not physics-forbidden.

The assumption stack is deep. Each item is individually plausible within current theoretical frameworks but experimentally unconfirmed. The conjunction of all six is speculative. The claim is not feasibility but self-consistency.

\subsection{Engineering versus physics gaps}
All internal physics questions identified during the design process have been resolved or quantitatively bounded: the Usov formula extrapolation, pair-capture efficiency, hadronization in strong magnetic fields, secondary-neutrino losses, sputtering yield, lateral stability and bulk emissivity through the firewall. The remaining open items are engineering: sustained 1000 T fields at 400 GPa magnetic pressure in the bore-throat solenoid, fabrication of 60,000 tonnes of CFL strange quark matter into a 3 pm spherical shell, construction of the 2,500 km SQM-bore linear accelerator and the bootstrap-recovery shell that must handle kT = 47 TeV during the seed phase. These requirements are extraordinary by any current standard. They are materials-science and systems-engineering challenges, not violations of known physical law. The distinction matters: a physics impossibility is permanent, while an engineering impossibility is contingent on the state of technology.

\subsection{The seed-production question}
The four-argument case against Drukier-Nussinov suppression at 10x threshold is the most speculative physics claim in this paper. The arguments --- classical restoration of GUT symmetry, the holy-grail function entering the over-barrier regime, Kibble freeze-out from a GUT-symmetric bubble and the vanishing WKB integral in the BPS limit --- are mutually reinforcing and draw on established sphaleron analogies [34, 32]. No lattice calculation exists for monopole production at finite collision energy in any specific GUT and non-perturbative corrections to the holy-grail function remain uncomputed. If Drukier-Nussinov suppression persists at all energies, the SQM linac fails and the monopole inventory must be obtained by other means. Both the committed path and the fallback are presented; the reader may judge the likelihood of each.

\subsection{Future work}
Five directions would most productively advance the design: (1) a full BlackHawk-PYTHIA simulation [35] of the Hawking spectrum at kT = 10.6 GeV; (2) lattice QCD calculations of the CFL gap at the operating chemical potential; (3) a non-perturbative evaluation of the holy-grail function for monopole production; (4) experimental monopole searches at MoEDAL [36] and future colliders; (5) a coupled black-hole-shell-ship dynamical simulation validating lateral stability beyond the linear regime.

\section{CONCLUSION}
The black hole propulsion concept [3] has been stalled for fifteen years by three unsolved problems: formation (kugelblitz formation are impossible [1]), energy capture (0.0001c limit [2]) and confinement (no mechanism proposed). This paper solves all three under six stated assumptions.

The key innovation is the CFL strange-quark-matter thermal wavelength converter: a shell that absorbs the full GeV Hawking spectrum via nuclear-density stopping and re-emits it as magnetically directable keV pairs through the Usov mechanism [10], raising capture from 46.5\% to 83\% and velocity from 0.0001c to 0.07---0.45c. The black hole is formed by BPS monopole collapse [29, 30] and grown from a 5.3 ms seed via railgun feeding. An anti-Helmholtz bottle confines the magnetically charged black hole and creates the exhaust nozzle.

The architecture offers three unique properties:

\begin{enumerate}
\item Thrust is continuously adjustable. Power scales as 1/M$^{2}$. At M = 10$^{9}$ kg: 55 MN, Alpha Centauri in 79 years at 0.07c. At M = 5 x 10$^{7}$ kg: 21,800 MN, 11 years at 0.45c.
\item The fuel is any matter. Conversion at mc$^{2}$ per kilogram regardless of composition.
\item Thermal management is solved by design. The Usov equilibrium is self-regulating: equilibrium temperature rises only logarithmically with input power. The three-layer outer firewall suppresses outward emission by 5 x 10$^{10}$ at every operating point. There is no radiator, no cooling loop and no thermal scaling limit on thrust. The thrust dial is real because the engine cannot overheat.
\end{enumerate}
The design is a thought experiment, not an engineering roadmap. The assumption stack is deep and untested. The physics chain is self-consistent, every numerical claim is derived from first principles with literature-anchored inputs and no known physical law is violated. The concept is conditional on its assumptions --- but it is no longer stuck on its open problems.

\section*{ACKNOWLEDGEMENTS}
This manuscript was prepared with AI assistance (Claude, Anthropic). All physics derivations were verified independently against the cited literature and reproduced numerically in companion calculation scripts. Detailed derivations and numerical calculation scripts are available from the corresponding author on request.

\section*{REFERENCES}
\small
\begin{enumerate}
\item A. Alvarez-Dominguez, L.J. Garay, E. Martin-Martinez and J. Polo-Gomez, 'No Black Holes from Light', Phys. Rev. Lett., 133, pp.041401, 2024.
\item J.S. Lee, 'Acceleration of a Schwarzschild Kugelblitz Starship', J. Brit. Interplanet. Soc., 68, pp.105, 2015.
\item L. Crane and S. Westmoreland, 'Are Black Hole Starships Possible?', arXiv:0908.1803, 2009.
\item K.F. Long, 'Deep Space Propulsion: A Roadmap to Interstellar Flight', Springer, New York, 2012.
\item S.W. Hawking, 'Particle Creation by Black Holes', Commun. Math. Phys., 43, pp.199-220, 1975.
\item D.N. Page, 'Particle Emission Rates from a Black Hole: Massless Particles from an Uncharged, Nonrotating Hole', Phys. Rev. D, 13, pp.198, 1976.
\item D.N. Page, 'Light Black Holes from Light', arXiv:2505.16202, 2025.
\item E.B. Bogomolny, 'Stability of Classical Solutions', Sov. J. Nucl. Phys., 24, pp.449, 1976.
\item M.K. Prasad and C.M. Sommerfield, 'Exact Classical Solution for the 't Hooft Monopole and the Julia-Zee Dyon', Phys. Rev. Lett., 35, pp.760, 1975.
\item V.V. Usov, 'Bare Quark Matter Surfaces of Strange Stars and $e^+e^-$ Emission', Phys. Rev. Lett., 80, pp.230, 1998.
\item J.P. Preskill, 'Magnetic Monopoles', Annu. Rev. Nucl. Part. Sci., 34, pp.461-530, 1984.
\item D. Olive and E. Witten, 'Supersymmetry Algebras That Include Topological Charges', Phys. Lett. B, 78, pp.97-101, 1978.
\item M. Alford, K. Rajagopal and F. Wilczek, 'Color-Flavor Locking and Chiral Symmetry Breaking in High Density QCD', Nucl. Phys. B, 537, pp.443-458, 1999.
\item M.G. Alford, A. Schmitt, K. Rajagopal and T. Schafer, 'Color Superconductivity in Dense Quark Matter', Rev. Mod. Phys., 80, pp.1455, 2008.
\item Y. Bai and T.-K. Chen, 'Excluding Stable Quark Matter: Insights from the QCD Vacuum Energy', arXiv:2502.20241, 2025.
\item O. Lennon, J. March-Russell, R. Petrossian-Byrne and S. Tillim, 'Black Hole Genesis of Dark Matter', J. Cosmol. Astropart. Phys., arXiv:1712.07664, 2018.
\item J. Maldacena, 'Comments on Magnetic Black Holes', J. High Energy Phys., 2021, pp.079, 2021.
\item T. Erber, 'High-Energy Electromagnetic Conversion Processes in Intense Magnetic Fields', Rev. Mod. Phys., 38, pp.626-659, 1966.
\item J.H. MacGibbon and B.R. Webber, 'Quark- and Gluon-Jet Emission from Primordial Black Holes: the Instantaneous Spectra', Phys. Rev. D, 41, pp.3052, 1990.
\item Particle Data Group, 'Review of Particle Physics', Phys. Rev. D, 110, pp.030001, 2024.
\item C. Manuel and K. Rajagopal, 'Illuminating Dense Quark Matter', Phys. Rev. Lett., 88, pp.042003, 2002.
\item V.V. Usov, 'Thermal Emission from Bare Quark Surfaces of Hot Strange Stars', Astrophys. J. Lett., 550, pp.L179, 2001.
\item M. Prakapenia and G.V. Vereshchagin, 'Revised Pair-Emission Rates from Bare Quark Surfaces of Strange Stars', Phys. Rev. D, 109, pp.023007, 2024.
\item P. Jaikumar, D.H. Rischke and I.A. Shovkovy, 'Goldstone-Boson Emissivities of Color-Superconducting Quark Matter', Phys. Rev. C, 68, pp.045803, 2003.
\item C. Alcock, E. Farhi and A. Olinto, 'Strange Stars', Astrophys. J., 310, pp.261, 1986.
\item M. Alford, C. Kouvaris and K. Rajagopal, 'Gapless Color Superconductors with Massed Quarks', Phys. Rev. D, 71, pp.054009, 2005.
\item G. 't Hooft, 'Magnetic Monopoles in Unified Gauge Theories', Nucl. Phys. B, 79, pp.276-284, 1974.
\item A.M. Polyakov, 'Particle Spectrum in Quantum Field Theory', JETP Lett., 20, pp.194-195, 1974.
\item N.S. Manton, 'The Force Between 't Hooft-Polyakov Monopoles', Nucl. Phys. B, 126, pp.525-541, 1977.
\item G.W. Gibbons and N.S. Manton, 'Classical and Quantum Dynamics of BPS Monopoles', Nucl. Phys. B, 274, pp.183-224, 1986.
\item A.K. Drukier and S. Nussinov, 'Monopole Pair Creation in Energetic Collisions: Is It Possible?', Phys. Rev. Lett., 49, pp.102, 1982.
\item A. Ringwald, 'High-Energy Breakdown of Perturbation Theory in the Electroweak Instanton Sector', Nucl. Phys. B, 330, pp.1-18, 1990.
\item T.W.B. Kibble, 'Topology of Cosmic Domains and Strings', J. Phys. A, 9, pp.1387-1398, 1976.
\item N.S. Manton, 'Topology in the Weinberg-Salam Theory', Phys. Rev. D, 28, pp.2019, 1983.
\item A. Arbey and J. Auffinger, 'BlackHawk: A Public Code for Calculating the Hawking Evaporation Spectra of any Black Hole Distribution', Eur. Phys. J. C, 79, pp.693, 2019.
\item B. Acharya et al. (MoEDAL Collaboration), 'MoEDAL Search in the CMS Beam Pipe for Magnetic Monopoles Produced via the Schwinger Effect', Phys. Rev. Lett., 133, pp.071803, 2024.
\item M. Ambrosio et al. (MACRO Collaboration), 'Final Results of Magnetic Monopole Searches with the MACRO Experiment', Eur. Phys. J. C, 25, pp.511-522, 2002.
\end{enumerate}
\normalsize
\end{document}